Let $n\in\mathbb{N}$ with $n\ge 1$.
\begin{itemize}
\item Define $h(n)$ to be the maximum integer $\ell\ge 0$ for which there exist \emph{distinct} primes
  \[p_1<p_2<\cdots<p_\ell\]
  such that $p_i\mid n$ for all $i$ and $p_{i+1}\equiv 1\pmod{p_i}$ for all $1\le i<\ell$.
  (If $n=1$, then $h(1)=0$.)
\item Define $H(n)$ to be the maximum integer $u\ge 1$ for which there exist integers
  \[1\le d_1<d_2<\cdots<d_u\]
  such that $d_i\mid n$ for all $i$ and $d_{i+1}\equiv 1\pmod{d_i}$ for all $1\le i<u$.
  (Note $1\mid n$, so $H(n)\ge 1$ always.)
\end{itemize}
Tasks:
\begin{enumerate}
\item Give meaningful upper/lower bounds (``estimates'') for $h(n)$ and $H(n)$ in terms of $n$ and standard arithmetic functions.
\item Decide whether $H(n)/h(n)\to\infty$ for \emph{almost all} $n$ (natural density $1$).
\end{enumerate}
